\documentclass[11pt]{article}

\usepackage[margin=1in]{geometry}
\usepackage{amsmath,amssymb,amsthm}
\usepackage{bm}
\usepackage{mathtools}
\usepackage{hyperref}
\usepackage{microtype}
\usepackage{booktabs}
\usepackage{xcolor}
\usepackage{listings}

\lstset{
  basicstyle=\ttfamily\small,
  breaklines=true,
  columns=fullflexible,
  showstringspaces=false,
  language=R,
}

\newcommand{\R}{\mathbb{R}}
\newcommand{\E}{\mathbb{E}}
\newcommand{\T}{^{\mathsf T}}
\newcommand{\inv}{^{-1}}
\DeclareMathOperator{\diag}{diag}

\title{Analytical Sequential State \& Disturbance Assimilation: \\
       The \texttt{SSSD.model} Function}
\author{}
\date{}

\begin{document}
\maketitle

\noindent
This note documents the data-assimilation routine
\texttt{SSSD.model} in
\texttt{Analytical\_OR\_original\_model.Rmd}, focusing on the
configuration \texttt{useAdjustment=TRUE} and
\texttt{use\_analytical\_posterior=TRUE}.  The filter is a
mixture-of-Gaussians extension of an Ensemble Adjustment Kalman
Filter (EAKF) in which each mixture component corresponds to a
distinct disturbance class.  At every annual timestep $t$ the
algorithm performs three steps:
\begin{enumerate}
\item a deterministic ecosystem \emph{forecast} with the VSEM model
      followed by stochastic process error and disturbance perturbation,
\item an \emph{analysis} step in which a closed-form Gaussian posterior
      is obtained per disturbance class and the class weights are
      updated by the marginal likelihood, and
\item an \emph{ensemble adjustment} (\texttt{multi.analytical.ens.adj})
      that deterministically rescales the forecast ensemble so its
      empirical moments and class frequencies match the analytical
      posterior, providing initial conditions for the next year.
\end{enumerate}

\paragraph{State and notation.}
Let $\bm x_t \in \R^{3}$ denote the carbon-pool state
$(C_v, C_s, C_w)$ -- foliage, soil and wood -- for a site at the end of
year $t$.  An ensemble of $n_e$ members
$\{\bm x_t^{(i)}\}_{i=1}^{n_e}$ approximates the filtering distribution.
Each member is also assigned a disturbance class
$c^{(i)} \in \{0, 1, \dots, n_d\}$, with $0 = $ undisturbed.  We write
$N = n_d + 1$ for the total number of mixture components and reserve
indices $d = 1, \dots, N$ for them ($d = 1$ corresponds to the
undisturbed component).

\bigskip

\section{Forecast: VSEM and the disturbance mixture}
\label{sec:forecast}

\subsection{The VSEM model}

VSEM (the ``Very Simple Ecosystem Model'' from the \texttt{BayesianTools}
R package) is a daily-time-step terrestrial carbon model with three
pools: foliage $C_v$, soil $C_s$ and wood $C_w$.  Given a vector of
$11$ ecosystem parameters $\bm\theta$ -- light-use efficiency, the
fraction of GPP allocated to foliage, root and wood, the turnover rates
of each pool, etc.\ -- and a daily PAR (photosynthetically active
radiation) driver, VSEM integrates the system
\[
  \frac{d}{dt}
  \begin{pmatrix} C_v \\ C_s \\ C_w \end{pmatrix}
  =
  f_{\mathrm{VSEM}}\!\bigl(C_v, C_s, C_w; \bm\theta,
                            \mathrm{PAR}(t)\bigr)
\]
forward by one day at a time.  In this code we use the daily increment
operator
\[
  \mathcal V_{\Delta t}\!\bigl(\bm x;\,\bm\theta,
                                \mathrm{PAR}\bigr)
\]
to push a state $\bm x$ forward over a window of length $\Delta t$ days.

\subsection{Per-member deterministic forecast}

At the start of timestep $t$ the analysis state from the previous year is
copied into entries $9{:}11$ of each ensemble member's parameter vector
(the ``initial pools'' slot of \texttt{THETA}):
\[
  \bm\theta^{(i)}_{9:11} \leftarrow \bm x^{(i),a}_{t-1}.
\]
VSEM is then run for one year (365 days) per member with that year's
PAR window, and the end-of-year pools are extracted:
\[
  \bm \mu^{(i)}_t
  \;=\;
  \mathcal V_{365}\!\bigl(\bm x^{(i),a}_{t-1};\,
                          \bm\theta^{(i)},
                          \mathrm{PAR}_{(t-1)\cdot 365 + 1 : t\cdot 365}
                    \bigr).
\]

\subsection{Adaptive process noise on $C_w$}

A Gaussian process-noise covariance $\bm Q \in \R^{3\times 3}$ is added
to the deterministic forecast.  The diagonal entry for the wood pool,
$Q_{ww}$, is updated sequentially via a conjugate Gamma update on the
\emph{precision} $\tau = 1/Q_{ww}$.  We assume
$Y_{t,w} \mid \bm \mu_t \sim
 \mathcal N\!\bigl(\bar\mu_{t,w},\, Q_{ww}\bigr)$
and place the prior
$\tau \sim \mathrm{Gamma}(a_0, b_0)$
with $a_0 = 1/0.1172475$, $b_0 = 1$.  After observing
$Y_{t,w}$ in a year that is neither a known disturbance year nor the
spin-up year, the conjugate update is
\[
  a_t \leftarrow a_t + \tfrac12,
  \qquad
  b_t \leftarrow b_t + \tfrac12 \bigl(Y_{t,w} - \bar\mu_{t,w}\bigr)^2,
\]
and we set
\[
  Q_{ww} \;\leftarrow\; \E\!\bigl[\tau\inv\bigr] \;=\; \frac{b_t}{a_t}.
\]

\subsection{Stochastic forecast and the disturbance mixture}

Process noise is sampled per member,
\(
  \bm x^{(i),f}_t \sim \mathcal N\!\bigl(\bm\mu^{(i)}_t,\,\bm Q\bigr).
\)
A disturbance class is then drawn for each member from a flat ensemble
prior
\(
  \pi_{\mathrm{ens}} = (0.5,\, 0.5/n_d, \dots, 0.5/n_d),
\)
and members assigned to class $d \ge 1$ are perturbed by
\texttt{disturbance.p}, which kills a fraction of foliage and wood
according to the class-$d$ removal distribution
$\mathcal N(\bm\mu^{0}_d, \bm V^{0}_d)$.

For each disturbance class $d \in \{1,\dots,n_d\}$, the empirical
moments of the perturbed sub-ensemble give the per-class \emph{forecast
prior}:
\[
  \bm\mu^{f}_{d+1}
  = \frac{1}{|S_d|}\sum_{i \in S_d} \bm x^{(i),f}_t,
  \qquad
  \bm C^{f}_{d+1}
  = \mathrm{Cov}\bigl(\{\bm x^{(i),f}_t\}_{i \in S_d}\bigr),
  \qquad
  \bm P^{f}_{d+1} = \bigl(\bm C^{f}_{d+1}\bigr)\inv,
\]
where $S_d$ is the set of members in class $d$.  If a sub-ensemble is
too small or its covariance is singular, the algorithm falls back to the
class's prior moments.  The undisturbed component $d=1$ is computed
analogously from the unperturbed members.  All states are then truncated
at zero by the helper \texttt{tobit}.

The result of the forecast step is the Gaussian-mixture forecast
distribution
\begin{equation}
  p^{f}(\bm x_t)
  \;=\;
  \sum_{d=1}^{N} \pi^{f}_d\,
                 \mathcal N\!\bigl(\bm x_t;\, \bm\mu^{f}_d,\, \bm C^{f}_d\bigr),
  \label{eq:forecast-mixture}
\end{equation}
together with the prior weights $\pi_d \equiv pD[d]$ supplied by the
external disturbance product for year $t$.

\bigskip

\section{Analysis: per-component closed-form Gaussian update}
\label{sec:analysis}

\subsection{Observation model}

A scalar observation $y_t = Y_{t,H}$ is available for one of the three
state components (in the Oregon application, the wood pool $C_w$, so
$H = 3$).  We encode the observed component with the row vector
$\bm H \in \R^{1 \times 3}$ that has a single $1$ in column $H$, so
$\bm H \bm x_t = x_{t,H}$.  Observation noise is heteroskedastic and
proportional-plus-intercept in the observed value:
\[
  R_H \;=\; m\, y_t + b,
  \qquad (m, b) = \texttt{Rparams},
\]
with the likelihood $y_t \mid \bm x_t \sim \mathcal N(\bm H \bm x_t, R_H)$.

\subsection{Posterior derivation, per component}

Conditional on disturbance class $d$, the prior is Gaussian
$\bm x_t \sim \mathcal N(\bm\mu^{f}_d, \bm C^{f}_d)$ with precision
$\bm P^{f}_d$.  The conditional posterior under the Gaussian likelihood
is also Gaussian by conjugacy.  Working in canonical (precision /
information) form, the posterior precision is the sum of the prior
precision and the contribution from the data,
\begin{equation}
  \bm P^{a}_d
  \;=\; \bm P^{f}_d + \bm H\T R_H\inv \bm H,
  \label{eq:precision-update}
\end{equation}
so that the posterior covariance is
\begin{equation}
  \boxed{\;
  \bm C^{a}_d \;=\;
   \bigl(\bm P^{f}_d + \bm H\T R_H\inv \bm H\bigr)\inv.
  \;}
  \label{eq:CA}
\end{equation}
The information-form (precision-weighted mean) update is
$\bm P^{a}_d \bm\mu^{a}_d
   = \bm P^{f}_d \bm\mu^{f}_d + \bm H\T R_H\inv y_t$, hence
\begin{equation}
  \boxed{\;
  \bm\mu^{a}_d \;=\;
   \bm C^{a}_d
   \bigl(\bm P^{f}_d \bm\mu^{f}_d
         + \bm H\T R_H\inv y_t\bigr).
  \;}
  \label{eq:muA}
\end{equation}
This pair of formulas is computed per disturbance class in the loop
\begin{lstlisting}
CA[d,,] <- solve(pf + t(mat_H) %*% ((1/R_H) * mat_H))
muA[d,] <- CA[d,,] %*% (pf %*% muf + (t(mat_H) * (1/R_H) * Y[t,H]))
\end{lstlisting}
which directly transcribes \eqref{eq:CA} and \eqref{eq:muA}.

\subsection{Mixture-weight update via the marginal likelihood}

Treating the disturbance class $D$ as a discrete latent variable with
prior $\Pr(D=d) = \pi_d$, Bayes' rule gives
\[
  \Pr(D = d \mid y_t)
  \;=\;
  \frac{p(y_t \mid D=d)\,\pi_d}{\sum_{d'} p(y_t \mid D=d')\,\pi_{d'}}.
\]
The component marginal likelihood is obtained by integrating out
$\bm x_t$ against the Gaussian prior,
\[
  p(y_t \mid D=d)
  \;=\;
  \int \mathcal N(y_t;\, \bm H \bm x_t, R_H)\,
       \mathcal N(\bm x_t;\, \bm\mu^{f}_d, \bm C^{f}_d)\, d\bm x_t
  \;=\;
  \mathcal N\!\bigl(y_t;\,
                    \bm H \bm\mu^{f}_d,\;
                    \bm H \bm C^{f}_d \bm H\T + R_H\bigr).
\]
This is the standard predictive distribution of a linear-Gaussian
state-space model.  In code, the unnormalized weights and their
normalization are
\begin{lstlisting}
lik   <- dnorm(Y[t,H], mean = mat_H %*% muf,
                       sd   = sqrt(mat_H %*% Cf %*% t(mat_H) + R_H))
wA[d] <- lik * pD[d]
...
wA <- wA / sum(wA)
\end{lstlisting}
yielding the posterior mixture weights $w^{a}_d = \Pr(D=d \mid y_t)$.

\subsection{Posterior summaries}

Given the per-component posteriors $(\bm\mu^{a}_d, \bm C^{a}_d)$ and
weights $w^{a}_d$, the algorithm reports
\begin{itemize}
  \item the overall posterior mean
        $\bm\mu^{a} = \sum_d w^{a}_d \bm\mu^{a}_d$,
  \item the disturbed-component conditional mean
        $\bm\mu^{a}_D = \sum_{d\ge 2} \frac{w^{a}_d}{\sum_{d'\ge 2}w^{a}_{d'}}
                          \bm\mu^{a}_d$,
  \item the undisturbed mean $\bm\mu^{a}_N = \bm\mu^{a}_1$,
  \item the total disturbance probability
        $\Pr(D \neq 0 \mid y_t) = \sum_{d\ge 2} w^{a}_d$,
  \item and per-component credible intervals from the most likely
        component, $\arg\max_d w^{a}_d$.
\end{itemize}

\bigskip

\section{Ensemble adjustment: \texttt{multi.analytical.ens.adj}}
\label{sec:adjustment}

The analysis step yields per-class posterior moments
$(\bm\mu^{a}_d, \bm C^{a}_d, w^{a}_d)$, but the next forecast step
needs an explicit ensemble of \emph{members}
$\{\bm x^{(i),a}_t\}$ that VSEM can be initialized from.
\texttt{multi.analytical.ens.adj} produces this ensemble by
deterministically transforming the forecast members so that, class by
class, the empirical moments of the ensemble match the analytical
posterior moments.  This is the same idea that underlies the EAKF, here
generalized to a Gaussian mixture.

\subsection{Class reassignment}

Let
\(
  f_d \;=\; \frac{1}{n_e}\sum_{i} \mathbf 1\{c^{(i)}_f = d\}
\)
be the empirical frequency of class $d$ in the forecast ensemble, and
let $w^{a}_d$ be the analysis weights.  We need the new class labels
$c^{(i)}_a$ to have empirical frequencies that approximate $w^{a}_d$.
The routine moves members \emph{out of} classes whose weight
\emph{decreased} ($w^{a}_d < f_d$) and \emph{into} classes whose weight
\emph{increased}, in proportion to the positive part of the gap.
Define
\[
  \Delta_d \;=\; \max\!\bigl(w^{a}_d - f_d,\; 0\bigr),
  \qquad
  q_d \;=\; \Delta_d \,\Big/\, \textstyle\sum_{d'} \Delta_{d'}.
\]
Then for each class $d$ with $w^{a}_d < f_d$, every member currently in
class $d$ is reassigned to a new class drawn from the categorical
distribution $q_{1:N}$.  Members in classes that grew or stayed the
same keep their forecast label.  The resulting labels $c^{(i)}_a$ have
expected frequencies equal to $w^{a}_d$ while changing as few labels
as necessary.

\subsection{Whitening: forecast $\to$ standard Gaussian}

Within each forecast class $d$, decompose the forecast covariance via
SVD,
\(
  \bm C^{f}_d = \bm V^{f}_d \diag(\bm L^{f}_d) (\bm V^{f}_d)\T.
\)
Each member is then mapped to a standardized vector
\begin{equation}
  \bm z^{(i)}
  \;=\;
  \diag(\bm L^{f}_d)^{-1/2} (\bm V^{f}_d)\T
  \bigl(\bm x^{(i),f}_t - \bm\mu^{f}_d\bigr),
  \qquad i \in S_d.
  \label{eq:whiten}
\end{equation}
By construction, if the within-class forecast empirical moments equal
$(\bm\mu^{f}_d, \bm C^{f}_d)$ exactly, then
$\{\bm z^{(i)}\}_{i \in S_d}$ has zero empirical mean and identity
empirical covariance.

\subsection{Recoloring: standard Gaussian $\to$ analysis posterior}

Now decompose each \emph{analysis} covariance,
\(
  \bm C^{a}_d = \bm V^{a}_d \diag(\bm L^{a}_d) (\bm V^{a}_d)\T,
\)
and recolor each member according to its \emph{new} class label
$c^{(i)}_a = d$:
\begin{equation}
  \bm x^{(i),a}_t
  \;=\;
  \bm V^{a}_d \, \diag(\bm L^{a}_d)^{1/2}\, \bm z^{(i)}
  + \bm\mu^{a}_d.
  \label{eq:color}
\end{equation}
The composition \eqref{eq:whiten}--\eqref{eq:color} is an affine map
that exactly preserves higher-order ensemble structure (skewness, joint
dependencies between members) while shifting the per-class mean to
$\bm\mu^{a}_d$ and the per-class covariance to $\bm C^{a}_d$.  In other
words: each forecast member is mapped \emph{deterministically} to an
analysis member by ``unrotating'' it using the forecast-class moments
and ``re-rotating'' it using the analysis-class moments, with a
possible class swap in between.  Because the transformation is
deterministic, no Monte-Carlo sampling noise is added by the analysis
step.  Pseudocode:
\begin{lstlisting}
for k in 1..N:
    SVD of forecast Pf[k]:           V_f, L_f
    for each forecast member i in class k:
        z[i] = (1/sqrt(L_f)) * V_f^T (Xf[i] - mu.f[k])

for k in 1..N:
    SVD of analysis Pp.Pa[k]:        V_a, L_a
    for each member i with new class k (cA[i] == k):
        X_a[i] = V_a * diag(sqrt(L_a)) * z[i] + Pp.mu.a[k]
\end{lstlisting}

\subsection{Why this yields a valid analysis ensemble}

If the per-class forecast empirical moments equal the assumed
$(\bm\mu^{f}_d, \bm C^{f}_d)$, the post-adjustment per-class empirical
moments equal $(\bm\mu^{a}_d, \bm C^{a}_d)$ exactly, and the empirical
class frequencies equal $w^{a}_d$.  Thus the adjusted ensemble is a
size-$n_e$ representation of the Gaussian-mixture posterior
\[
  p^{a}(\bm x_t \mid y_t)
  \;=\;
  \sum_{d=1}^{N} w^{a}_d\,
                 \mathcal N\!\bigl(\bm x_t;\, \bm\mu^{a}_d,\, \bm C^{a}_d\bigr)
\]
that we derived in Section~\ref{sec:analysis}, and it can be fed
straight into the next year's VSEM forecast as the new ICs
$\bm x^{(i),a}_t$.  Negative pools that arise from the affine map are
truncated to zero by \texttt{tobit} before the next forecast.

\bigskip

\noindent
\textbf{Summary.}  Each annual cycle of \texttt{SSSD.model} (with the
analytical posterior and ensemble adjustment turned on) is a
mixture-of-Gaussians Kalman update: VSEM provides the per-class forecast
moments, the closed-form expressions \eqref{eq:CA}--\eqref{eq:muA} give
the per-class analysis moments, the marginal likelihood updates the
class weights, and \texttt{multi.analytical.ens.adj} converts those
moments back into a discrete ensemble for the next forecast step.

\end{document}
